\documentclass{article}
\usepackage[utf8]{inputenc}
\usepackage[a4paper, margin=3cm]{geometry}
\setlength{\parindent}{0pt} %--don't indent paragraphs
\begin{document}
\centerline{\Large\bfseries Capital Gains Tax Calculations for 2012-13}

\section*{Capital Gains Events}


\clearpage
\section*{Interest and Dividends Events}


\subsection*{01 June 2012}
\textbf{Dividend 1:} BAR for £100.00
(tax paid at source: £15.00)
, tax treaty amount: £15.00 (treaty rate for USA: 15.00\%)



\clearpage
\section*{Overall - Capital Gains}
\subsection*{Number of acquisitions: 0}
\subsection*{Number of disposals: 0}
\subsection*{Total disposal proceeds: £0.00}
\subsection*{Total capital gain before loss: £0}
\subsection*{Total capital loss: £0}
\subsection*{Total capital gain after loss: £0}
\vspace{1em}
Amounts are rounded to the nearest penny individually, so a line may differ from its components by a penny.
\vspace{2em}
\section*{Overall - Dividends and Interests}
\subsection*{Number of dividend distributions: 1}
\subsection*{Number of interest transactions: 0}
\subsection*{Number of interest tax transactions: 0}
\subsection*{Number of excess reported income distributions: 0}
\subsection*{Total amount of dividends proceeds: £100.00}
\subsection*{Total amount of dividends tax allowance due to double taxation treaties: £15.00}
Most dividends before 6 April 2016 carried a tax credit, so the taxable amount is not worked out.
\subsection*{Total amount of UK interests proceeds: £0.00}
\subsection*{Total amount of foreign interests proceeds: £0.00}
\subsection*{Total amount of interest tax paid: £0.00}
\end{document}